\chapter{\IfLanguageName{dutch}{Methodologie}{Methodology}}
\label{ch:methodologie}

Dit hoofdstuk geeft een overzicht van het plan van aanpak voor de uitvoering van het onderzoek. Het bouwt voort op de inzichten uit Hoofdstuk \ref{ch:stand-van-zaken} en dient als basis voor de uitvoering van de verschillende onderzoekstaken.

\section{Opstellen van een dataset}
\label{sec:opstellen-van-een-dataset}

Het opstellen van een dataset is essentieel voor het trainen en evalueren van het model voor het detecteren van journalistieke partijdigheid. Zoals beschreven in Hoofdstuk \ref{ch:stand-van-zaken} wordt een trainingsstrategie toegepast waarbij het model wordt getraind op drie verschillende datasets, die elk een specifieke rol vervullen binnen het classificatieproces.

\subsection{Metadata-dataset}

De metadata-dataset wordt verkregen door het scrapen van nieuwsartikelen van de \gls{hln}-website. Voor elk artikel worden relevante metadata verzameld, zoals:

\begin{itemize}
    \item Auteur
    \item Publicatiedatum
    \item Nieuwssectie
\end{itemize}

De scraping wordt uitgevoerd met Python en de \textit{BeautifulSoup}-bibliotheek, waarbij de \gls{dom}-structuur van de webpagina wordt geanalyseerd om de benodigde informatie te extraheren.

\subsection{Sentence-level dataset}

Voor training op zinsniveau wordt gebruikgemaakt van de \textit{FactNews}-dataset \autocite{Vargas2024}. Deze dataset wordt vertaald naar het Nederlands en bevat zinnen die gelabeld zijn op basis van partijdigheid. Hierdoor kan het model signalen van bias binnen afzonderlijke zinnen leren herkennen.

\subsection{Document-level dataset}

Voor het trainen op documentniveau wordt de \textit{NewsMediaBias-Plus}-dataset van \textcite{Raza2024} gebruikt. Ook deze dataset wordt vertaald naar het Nederlands en bevat labels op het niveau van volledige artikelen. Training op documentniveau helpt het model om bredere patronen van partijdigheid te herkennen die niet zichtbaar zijn op zinsniveau.

\subsection{Testset}

De testset bestaat uit artikelen van \gls{hln}, die gescraped worden met dezelfde aanpak als voor de metadata-dataset. Voor elk artikel worden de volgende elementen verzameld:

\begin{itemize}
    \item Titel van het artikel
    \item Volledige inhoud van het artikel
    \item Metadata (auteur, publicatiedatum, nieuwssectie)
\end{itemize}

De testset wordt gebruikt om de prestaties van het getrainde model te evalueren op echte \gls{hln}-artikelen, waarbij de voorspellingen van het model worden vergeleken met de labels van externe mediabias-tools.

\section{Opstellen van een pipeline}
\label{sec:opstellen-van-een-pipeline}

Voor de verwerking van de drie datasets wordt een pipeline ontwikkeld in Python. De pipeline bestaat uit drie componenten die elk afzonderlijk worden getraind op hun respectievelijke dataset. De outputs van de metadata- en sentence-level componenten worden geïntegreerd als aanvullende features binnen het document-level component, dat een eindvoorspelling genereert over de partijdigheid van het volledige artikel.

De pipeline bestaat uit de volgende fasen:

\begin{enumerate}
    \item Data-inname
    \item Preprocessing
    \item Feature-engineering
    \item Modeltraining en validatie
\end{enumerate}

De outputs van de metadata- en sentence-level pipelines geven invloed binnen de document-level pipeline.

\subsection{Preprocessing}

Tijdens de preprocessingfase worden de ruwe zinnen en artikelen omgezet naar een vorm die beter geschikt is voor het trainen van modellen. De volgende preprocessingstappen worden uitgevoerd in de aangegeven volgorde:

\begin{enumerate}
    \item Het verwijderen van HTML-tags en andere niet-tekstuele elementen
    \item Het vertalen van de tekst naar het Nederlands (indien nodig)
    \item Het omzetten van alle tekst naar kleine letters (\textit{lowercase conversion})
    \item Het splitsen van de tekst in afzonderlijke tokens (\textit{tokenization})
    \item Het verwijderen van stopwoorden (\textit{stop-word removal})
    \item Het herleiden van woorden naar hun correcte woordenboekvorm (\textit{lemmatization})
\end{enumerate}

\subsection{Feature-engineering}

Na de preprocessing wordt de tekst omgezet naar een numerieke representatie met behulp van contextuele woordembeddings. Daarnaast worden aanvullende features toegevoegd, zoals sentiment-scores en metadata-gebaseerde kenmerken, die kunnen bijdragen aan het detecteren van partijdigheid.

\subsection{Datasplit}

Voor het trainen en evalueren van de modellen worden de datasets opgesplitst in afzonderlijke subsets.

\begin{itemize}
    \item De trainingsdata wordt gebruikt om de modellen te trainen.
    \item De validatiedata wordt gebruikt om de prestaties tijdens het trainen te monitoren.
    \item De testdata wordt gebruikt voor de finale evaluatie.
\end{itemize}

Bij de splitsing wordt \textit{stratified sampling} toegepast om ervoor te zorgen dat de verhouding tussen de verschillende klassen behouden blijft in elke subset. Daarnaast wordt \textit{cross-validation} toegepast om de prestaties van het model te kunnen beoordelen op een gegeneraliseerde manier.

\section{Trainen en evalueren}
\label{sec:trainen-en-evalueren}

\subsection{Standaard evaluatiemetrieken}

De prestaties van de modellen worden geëvalueerd met behulp van standaard classificatiemetrieken, waaronder \textit{accuracy}, \textit{precision}, \textit{recall} en de \textit{F1-score}. Deze metrieken geven inzicht in hoe goed het model correcte voorspellingen maakt en hoe het omgaat met fouten.

\subsection{MBIB benchmark}

Naast standaardmetrieken wordt het model geëvalueerd met behulp van de \gls{mbib}-benchmark \autocite{Wessel2023}. Deze benchmark combineert meerdere datasets en maakt het mogelijk om de generalisatiecapaciteit van het model te evalueren over verschillende vormen van bias.

\subsection{Vergelijking met bestaande tools}

De resultaten van het model worden ook vergeleken met bestaande mediabias-tools, zoals \textit{Ground News} en \textit{Media Bias}. Hierbij wordt nagegaan in welke mate de voorspellingen overeenkomen en waar verschillen optreden.

\section{Implementatie van een browserextensie}
\label{sec:implementatie-van-een-browserextensie}

Een browserextensie wordt ontwikkeld in JavaScript om de partijdigheid van nieuwsartikelen automatisch te analyseren. De extensie verwerkt de DOM-structuur van een webpagina, extraheert de inhoud van het artikel en stuurt deze door naar het getrainde model.

De output van het model wordt visueel weergegeven binnen de webpagina. Daarnaast worden de zinnen die het meest bijdragen aan de voorspelling gemarkeerd in de tekst, zodat de gebruiker inzicht krijgt in de motivatie van de classificatie.

\section{Resultatenanalyse}
\label{sec:resultatenanalyse}

De resultaten van de evaluatiefase en benchmarktesten worden geanalyseerd, waarbij gekeken wordt naar:

\begin{itemize}
    \item De prestaties van de individuele componenten.
    \item De prestaties van de volledige pipeline op de testset van \gls{hln}-artikelen.
    \item Een vergelijking van de voorspellingen van het model met die van bestaande mediabias-tools.
\end{itemize}

Naast algemene evaluatiemetrieken wordt gebruikgemaakt van een \textit{confusion matrix}, zie Tabel \ref{tab:confusion-matrix}, om de verschillende types fouten te analyseren.

\begin{table}
  \centering
  \begin{tabular}{c|c|c}
     & \textbf{Voorspeld: Neutraal} & \textbf{Voorspeld: Partijdig} \\
    \hline
    \textbf{Werkelijk: Neutraal} & True Negative (TN) & False Positive (FP) \\
    \hline
    \textbf{Werkelijk: Partijdig} & False Negative (FN) & True Positive (TP) \\
  \end{tabular}
  \caption[Confusion matrix]{Overzicht van correcte en foutieve classificaties van het model.}
  \label{tab:confusion-matrix}
\end{table}

De confusion matrix maakt het mogelijk om inzicht te krijgen in het aantal \textit{false positives} en \textit{false negatives}.